../cictikz/src/cictikz/data/tex/cictikz_preamble.tex

% Curvature correction for the current-mode (Banba) reference of l3_ptat3.
%
% Two identical PNPs are biased at currents with different temperature
% exponents: Q_A from a PTAT current, Q_B from the compensated current the core
% already makes. Their V_BE difference is exactly (kT/q) ln(I_A/I_B), and since
% that ratio is itself proportional to T, the difference carries a
% (kT/q) ln(T/T0) term -- the T ln T shape the residual curvature needs.
%
% The OTA holds the right end of R4 at V_BE,A while its left end sits at
% V_BE,B, so I_NL = (V_BE,A - V_BE,B)/R4. MP_C is the pass device in that loop
% and MP_D copies I_NL into the R3 summing node from Figure 15.
%
% Feedback polarity: raising the OTA output turns MP_C off, which lowers the
% node it feeds. The sensed node is therefore the + input and the fixed V_BE,A
% is the - input, or the loop latches instead of settling.

\begin{circuitikz}[american, thick, transform shape, circuit ee IEC]

  ../cictikz/src/cictikz/data/tex/cictikz_lib.tex

  \def\xa{0}             % Q_A, biased by a PTAT current
  \def\xb{\grid*2.5}     % Q_B, biased by the compensated current
  \def\xn{\grid*6}       % MP_C, pass device of the R4 loop
  \def\xd{\grid*8.5}     % MP_D, the copy into the output node
  \def\ynl{4.2}          % R4 / node X level, also the PMOS drain level
  \def\ypds{5.8}         % PMOS source level
  \def\ytop{5.4}         % bottom of the two bias current sources
  \def\ysup{7.0}         % supply rail
  \def\yo{3.4}           % OTA + input

  % ---- Q_A: diode connected PNP at a PTAT current ----
  \node[pnp, anchor=C, label=right:$Q_{A}$] (QA) at (\xa,0) {};
  \draw (QA.C) -- ++(0,-0.2) \vground;
  \draw (QA.B) |- (QA.C);
  \draw (QA.E) -- (\xa,\ytop);
  \draw (\xa,\ysup) to [I, l=$I_{PTAT}$] (\xa,\ytop);
  \draw (\xa,\ysup) \vsupply;

  % ---- Q_B: same device, biased by the temperature compensated current ----
  \node[pnp, anchor=C, label=right:$Q_{B}$] (QB) at (\xb,0) {};
  \draw (QB.C) -- ++(0,-0.2) \vground;
  \draw (QB.B) |- (QB.C);
  \draw (QB.E) -- (\xb,\ytop);
  \draw (\xb,\ysup) to [I, l=$I_{REF}$] (\xb,\ytop);
  \draw (\xb,\ysup) \vsupply;

  % ---- R4 from node b out to node X ----
  \draw (\xb,\ynl) \hresistor{$R_{4}$};
  \coordinate (r4right) at (resEnd);
  \draw (\xn,\ynl) to [short, i=$I_{NL}$] (r4right);
  \fill (\xb,\ynl) circle (0.06);

  % ---- OTA: + senses node X, - sits on V_BE,A ----
  \draw (\xa+\grid*4,\yo) \cicOtaSSWP{$+$}{$-$};
  \draw (cicOtaS_inp) -- (\xb+1.8,\yo) -- (\xb+1.8,\ynl);
  \fill (\xb+1.8,\ynl) circle (0.06);
  % Crosses the Q_B branch on the way to node a -- no dot there.
  \draw (cicOtaS_inn) -- (\xa,\yo-0.8);
  \fill (\xa,\yo-0.8) circle (0.06);

  % ---- Pass device and the mirrored copy ----
  \draw (\xn,\ynl) \lvmpmos{MPC}{};
  \draw (MPC.source) \vsupply;
  \draw (\xd,\ynl) \lvpmos{MPD}{};
  \draw (MPD.source) \vsupply;
  \draw (MPC.gate) -- (MPD.gate);

  \draw (cicOtaS_out) -- (\xn+1.8,\yo-0.4);
  \draw (\xn+1.8,\yo-0.4) |- (MPC.gate);
  \fill (\xn+1.8,5.0) circle (0.06);

  % ---- Output: I_NL lands on the summing node of Figure 15 ----
  \draw (\xd,\ynl) to [short, i=$I_{NL}$] (\xd,2.2);
  \draw (\xd,0) \vground;
  \draw (\xd,0) -- ++(0,0.4);
  \draw (\xd,0.4) \vresistor{$R_{3}$};
  \draw (resEnd) -- (\xd,2.2);
  \draw (\xd,2.2) \portOut{$V_{REF}$};
  \draw (\xd,2.2) to [short,*-o] (\xd-1.4,2.2);
  \node[anchor=east] at (\xd-1.5,2.2) {$I_{PTAT}+I_{CTAT}$};

\end{circuitikz}
\end{document}
